\usepackage{etoolbox} % Permite lógica booleana e condicionais
\usepackage{pgffor}   % Permite loops

% --- DEFINIÇÃO DAS VARIÁVEIS (GETTERS/SETTERS) ---
%> Variável booleana (true/false) para formatação
\newtoggle{ufvnorms}
\newcommand{\UseUFVNorms}[1]{%
    \ifstrequal{#1}{True}{\toggletrue{ufvnorms}}{%
    \ifstrequal{#1}{true}{\toggletrue{ufvnorms}}{\togglefalse{ufvnorms}}}%
}
%> Variável boolena (true/false) para o cabeçalho
\newtoggle{useheaders}
\newcommand{\UseHeaders}[1]{%
    \ifstrequal{#1}{True}{\toggletrue{useheaders}}{%
    \ifstrequal{#1}{true}{\toggletrue{useheaders}}{\togglefalse{useheaders}}}%
}
\newtoggle{numericformat}
\toggletrue{numericformat}
\newcommand{\NumericFormat}[1]{%
    \ifstrequal{#1}{True}{\toggletrue{numericformat}}{%
    \ifstrequal{#1}{true}{\toggletrue{numericformat}}{\togglefalse{numericformat}}}%
}
\newtoggle{pretextual}
\toggletrue{pretextual}
\newcommand{\PreTextual}[1]{%
    \ifstrequal{#1}{True}{\toggletrue{pretextual}}{%
    \ifstrequal{#1}{true}{\toggletrue{pretextual}}{\togglefalse{pretextual}}}%
}
\NewDocumentCommand{\SetMathematicalFont}{ O{TeX Gyre Pagella Math} }{
    \def\MathFont{#1}
}
\NewDocumentCommand{\SetTextFont}{ O{TeX Gyre Pagella} }{
    \def\TextFont{#1}
}

%> Comandos para nomear o documento
\NewDocumentCommand{\DocTitle}{ O{Título não definido}}{
    \def\doctitle{#1}}
\NewDocumentCommand{\DocAuthor}{ O{Autor não definido} }{
    \def\docauthor{#1}}
\NewDocumentCommand{\DocAdvisor}{ O{} }{
    \def\docadvisor{#1}}
\NewDocumentCommand{\DocCoAdvisor}{ O{} }{
    \def\doccoadvisor{#1}}
\NewDocumentCommand{\DocModality}{ O{Dissertação} }{
    \def\docmodality{#1}}
\NewDocumentCommand{\DocDegree}{ O{Mestre} }{
    \def\docdegree{#1}}
\NewDocumentCommand{\DocDate}{ O{\today} }{
    \def\docdate{#1}}
\NewDocumentCommand{\DocProgram}{ O{Física} }{
    \def\docprogram{#1}}
\NewDocumentCommand{\DocCenter}{ O{Centro de Ciências Exatas} }{
    \def\doccenter{#1}}
\NewDocumentCommand{\DocPresentation}{ O{} }{
    \def\docpresentation{#1}}
\NewDocumentCommand{\ReportCourse}{ O{} }{
    \def\reportcourse{#1}}
\NewDocumentCommand{\DocAbstract}{ O{} }{
    \def\docabstract{#1}}
\NewDocumentCommand{\DocResumo}{ O{} }{
    \def\docresumo{#1}}
\NewDocumentCommand{\DocKeywords}{ O{} }{
    \def\dockeywords{#1}}
\NewDocumentCommand{\DocPalavrasChave}{ O{} }{
    \def\docpalavraschave{#1}}
\NewDocumentCommand{\DocDedication}{ O{} }{
    \def\docdedication{#1}}
\NewDocumentCommand{\DocAcknowledgments}{ O{} }{
    \def\docacknowledgments{#1}}
\NewDocumentCommand{\DocEpigraph}{ O{} O{} }{
    \def\docepigraph{#1}%
    \def\epigraphauthor{#2}}
\NewDocumentCommand{\SetStartingPage}{ O{} }{
    \def\docstartingpage{#1}}
\NewDocumentCommand{\BoardNames}{ O{} }{
    \def\boardnames{#1}}

%> Comandos para nomear o documento
\SetMathematicalFont
\SetTextFont
\DocTitle
\DocAuthor
\DocAdvisor
\DocCoAdvisor
\DocModality
\DocDegree
\DocDate
\DocProgram
\DocCenter
\DocPresentation
\ReportCourse
\DocAbstract
\DocResumo
\DocKeywords
\DocPalavrasChave
\DocDedication
\DocAcknowledgments
\DocEpigraph
\SetStartingPage
\BoardNames

%> Macro para construção do template
\newcommand{\BuildTemplateSetup}{%
    
    %> 1. Carrega os pacotes universais
    % --- PACOTES ---
%> Pacotes essenciais
\usepackage{graphicx}
\usepackage{indentfirst}
\usepackage[utf8]{inputenc}
\usepackage[T1]{fontenc}
\usepackage[brazilian]{babel}

%> Pacotes de estilo
\usepackage{titlesec}
\usepackage{multicol}
\usepackage{epigraph}
\usepackage{capt-of}
\usepackage{microtype}
\usepackage[autostyle]{csquotes}
\usepackage{booktabs}
\usepackage{enumitem}
\usepackage{array}
\renewcommand{\arraystretch}{1.6}
\usepackage[svgnames]{xcolor}
\usepackage{pdfpages} %> Necessário para adicionar páginas em pdf (registro, etc...)
\usepackage{fancyhdr}
\setlength{\headheight}{14.5pt}
%\usepackage{markdown} %> Para escrever linhas de código

%> Pacotes matemáticos
\usepackage{amsmath, amssymb, amsfonts}
\usepackage{esint}
\usepackage{upgreek}
\usepackage{tikz}    
\usetikzlibrary{angles, quotes, arrows.meta, babel}
\usepackage{xsim}
\usepackage[most]{tcolorbox}

%> Configurações das páginas
\pagestyle{plain}

%> Gerador de blá blá blá
\newcommand{\blablabla}[1][1]{%
  \foreach \n in {1,...,#1}{%
    Neste parágrafo vamos simular um texto em português para testar a tipografia, o espaçamento, a hifenização e a formatação geral do documento. Como estamos desenvolvendo um modelo robusto para a Universidade Federal de Viçosa (UFV), é importante que as margens sejam testadas com palavras reais da nossa língua, contendo acentuações, cedilhas e construções frasais típicas do nosso idioma. Dessa forma, garantimos que o código do template funcione perfeitamente, sem depender de pacotes externos ou línguas mortas. \par\vspace{0.2cm}%
  }%
}
    \usepackage{unicode-math}
    \setmathfont{\MathFont}
    \usepackage{hyperref}

    %> 2. Carrega as opções de formatação: UFV / Publish
    \iftoggle{ufvnorms}{%
        \usepackage[a4paper,
            left=3cm, top=3cm, right=2cm, bottom=2cm]{geometry}
        \input{_preamble/format2ufv}     % Carrega o driver da UFV
    }{%
        \usepackage[a4paper,
            margin=2.5cm, top=2.5cm, bottom=2.5cm]{geometry}
        \setmainfont{\TextFont}
        \usepackage[a4paper,
        margin=2.5cm,
        top=2.5cm,
        bottom=2.5cm
    ]{geometry}
\usepackage[
    backend=biber,
    style=numeric-comp,
    sorting=none,
    natbib=true
]{biblatex}

\iftoggle{useheaders}{
    \pagestyle{fancy}
    \fancyhf{}
    \fancyhead[R]{\thepage}
    \fancyhead[L]{\slshape\nouppercase{\leftmark}} % Itálico (slanted)
    \renewcommand{\headrulewidth}{0.4pt} % Linha charmosa separando o header
}{
    \pagestyle{plain} 
}

\usepackage{mathpazo}

%> Formatação dos capítulos enumerados
\titleformat{\chapter}[display]
    {\normalfont\bfseries}
    {\filleft\LARGE\chaptertitlename\ \thechapter} % Alinhamento à direita para "Capítulo X"
    {0.5cm} % Espaço entre o número e a linha
    {\titlerule\vspace{0.8cm}\filright\Huge} % Título do capítulo alinhado à esquerda
    [\vspace{2cm}]
  
\titleformat{name=\chapter,numberless}[display]
  {\normalfont\bfseries}{}{0pt}
  {\filright\Huge} % AQUI REMOVEMOS O \titlerule
  [\vspace{2cm}]

\titlespacing*{\chapter}{0pt}{0pt}{0pt}

%> Configuração das epígrafes dos capítulos
\renewcommand{\epigraph}[2]{
  \begin{flushright}
  \begin{minipage}{.3\textwidth}
    \hspace{0.3cm}{#1}\vspace{0.3cm}\\
    {\itshape\small#2}
  \end{minipage}\hspace*{2cm}
\end{flushright}
\vspace{1.5cm}
} % Carrega o driver fancy
    }

    %> 3. Configura a notação padrão, o hyperref e o bibtex
    %% --NOTAÇÃO PADRÃO--

%> Alfabetos matemáticos e operadores
\setmathfont[range={scr, bfscr}]{XITS Math}
\let\dcal\mathscr
\setmathfont[range={cal, bfcal}, Scale=MatchUppercase]{Asana Math}
\let\pzcal\mathcal

%> Notação para variáveis e grandezas
%> Diferencial
\newcommand{\diff}{
    \mathop{}\!\symup{d}}
    \iftoggle{ufvnorms}{%
        \hypersetup{
            colorlinks=true, 
            allcolors=black
        }
    }{%
        \hypersetup{
            colorlinks=true,
        	linkcolor=NavyBlue,
        	urlcolor=MediumBlue,
        	citecolor=ForestGreen,
        	pdfhighlight=/N
        }

    }
    \addbibresource{_bibliography/references.bib}

    %> 4. Metadados no PDF e no Título do LaTeX
    \hypersetup{
        pdftitle={\doctitle},
        pdfauthor={\docauthor}
    }
    \title{\doctitle}
    \author{\docauthor}
    \usepackage{cleveref}

    %> 5. Carregar os macros pré-textuais
    \newcommand{\Cabecalho}{
    Universidade Federal de Viçosa\\
    \thesiscenter\\ % Centro...
    % Lógica condicional: Departamento (Graduação) vs Programa (Pós)
    \ifdefstring{\thesismodality}{Monografia}{%
    Departamento de \thesisprogram\\
    }{%
    Programa de Pós-Graduação em \thesisprogram\\
    }
}

\newcommand{\BigTitle}{
    \begin{center}
    \rule{0.4\linewidth}{0.8pt}\vspace{0.4cm}\par
    {\bfseries\Large \thesistitle \par}
    \vspace{0.3cm}\rule{0.4\linewidth}{0.8pt}\par
    \end{center}
}

\newcommand{\bfTitle}{
    {\bfseries\thesistitle}
}

\newcommand{\DocPresentation}{
    \hfill
    \begin{minipage}[c][3cm][c]{8cm}
    \ifdefstring{\thesismodality}{Monografia}{%
        \small{\thesismodality\ apresentada ao Departamento de \thesisprogram\ da Universidade Federal de Viçosa como parte das exigências para a conclusão da disciplina \reportcourse.}
    }{%
        \small{\thesismodality\ apresentada à Universidade Federal de Viçosa como parte das exigências para obtenção do título de \thesisdegree\ em \thesisprogram.}
    }
    \end{minipage}
}

\newcommand{\DocNames}{
    \hfill
    \begin{minipage}[c][2cm][c]{8cm}
    \small{
    \begin{flushright}
    \textbf{Estudante:} \thesisauthor \\
    \textbf{Orientador:} \thesisadvisor \\
    \ifdefempty{\thesiscoadvisor}{}{%
        \textbf{Coorientador:} \thesiscoadvisor%
    }
    \end{flushright}
    }
    \end{minipage}
}

\newcommand{\DocLocation}{
    \begin{center}
    Viçosa, Minas Gerais -- Brasil\\
    \thesisdate
    \end{center}
}

\newcommand{\BuildAcknowledgments}{
    \begin{center}
        {\bfseries\LARGE Agradecimentos}\vspace{0.5cm}
    \end{center}

    \hrulefill
    \vspace{2cm}

    \docacknowledgments
}

\newcommand{\BuildEpigraph}{
    \vspace*{\fill}
    \begin{flushright}
    \begin{minipage}{0.5\textwidth}
        \begin{flushleft}
                {\itshape
                \docepigraph}
        \end{flushleft}
        \begin{flushright}
                --- \epigraphauthor
        \end{flushright}
    \end{minipage}
    \end{flushright}
    \vspace{2cm}
}

\newcommand{\BuildResumo}{
    \begin{center}
    \begin{minipage}[c][3cm][c]{12cm}
    \centering
    \Cabecalho
    \vspace{0.3cm}
    \textbf{Resumo}\\
    \vspace{0.3cm}
    \bfTitle
    \vspace{1cm}
    \end{minipage}
    \end{center}
    
    \docresumo\\

    \noindent
    \textbf{Palavras-chave:} \docpalavraschave.
    
}

\newcommand{\BuildAbstract}{
    \begin{center}
    \begin{minipage}[c][3cm][c]{12cm}
    \centering
    \Cabecalho
    \vspace{0.3cm}
    \textbf{Abstract}\\
    \vspace{0.3cm}
    \bfTitle
    \vspace{1cm}
    \end{minipage}
    \end{center}
    
    \docabstract\\

    \noindent
    \textbf{Keywords:} \dockeywords.
    
}


\newcommand{\BuildPretextual}{
\iftoggle{pretextual}{
    
    \pagenumbering{Alph}
    \newcommand{\BuildTitlePage}{
\begin{titlepage}
\noindent\begin{minipage}[t][6cm][t]{\textwidth}
    \begin{center}
        \Cabecalho\vspace{2cm}

        {\large \textbf{\docauthor}}
    \end{center}
\end{minipage}

\vspace*{\fill}
\BigTitle
\vspace*{\fill}

\noindent\begin{minipage}[b][6cm][b]{\textwidth}
    \DocLocation
\end{minipage}
\end{titlepage}
}

\newcommand{\BuildCoverPage}{
\cleardoublepage
\begin{titlepage}
\noindent\begin{minipage}[t][9cm][t]{\textwidth}
    \begin{center}
        {\large \textbf{\docauthor}}
    \end{center}
\end{minipage}

\vspace*{\fill}
\BigTitle
\vspace*{\fill}

\noindent\begin{minipage}[b][9cm][b]{\textwidth}
    \BuildPresentation
    \vspace{0.5cm}
    \DocNames
    \vspace*{\fill}
    \DocLocation
\end{minipage}
\end{titlepage}
}

\ifboolexpr{
    test {\ifdefstring{\docmodality}{Monografia}} 
    or 
    test {\ifdefstring{\docmodality}{Relatório}}
}{
    % SE FOR MONOGRAFIA OU RELATÓRIO:
    \BuildCoverPage
}{
    % DISSERTAÇÃO OU TESE:
    \BuildTitlePage
    \BuildCoverPage
}

    \cleardoublepage
    \pagenumbering{roman}

    \ifdefempty{\docresumo}{}{
    \newpage
    \BuildResumo    }

    \ifdefempty{\docabstract}{}{
    \newpage        
    \BuildAbstract  }

    \ifdefempty{\docacknowledgments}{}{
    \newpage
    \BuildAcknowledgments   }

    \ifdefempty{\docepigraph}{}{
    \newpage
    \BuildEpigraph  }
    }{}
\newpage
\tableofcontents
\cleardoublepage
\pagenumbering{arabic}
}
}